



% Jan.  28, 2004. Submitted to the Proceedings of Hayama Conference, 2003
\input amstex
\documentstyle{amsppt}\magnification=\magstep1
\magnification=1200
%\hcorrection{0.25in}
%\vcorrection{-0.8in}
\NoRunningHeads
\NoBlackBoxes
\baselineskip=12pt
%\parindent=10pt
\vsize=8.5truein
\hsize=6.0truein
\hoffset=0.3truein
\voffset=0.2truein
\loadmsam
\loadbold
\loadmsbm
\loadeufm \loadeusm \UseAMSsymbols
%\font\boldtitlefont=cmb10 scaled\magstep2 \font\sectionfont=cmb10
%scaled\magstep1
     \NoRunningHeads
     \nologo
\document
\define\({\left(}
\define\){\right)} \define\[{\left[} \define\]{\right]}
\define\db{\bar\partial} \define\p{\parallel}
\define\aint{\rlap{\kern2pt\vrule height2.5pt width9.5pt depth-2.1pt}
\int_{B_{\rho}}}
       \nologo

\topmatter
\title Estimates for the $\db$-Neumann problem     and nonexistence of
$C^2$
     Levi-flat hypersurfaces  in $\Bbb CP^n$*
      \endtitle

\author   Jianguo Cao,     Mei-Chi Shaw  and Lihe Wang  \endauthor

%Department of Mathematics, University of Notre Dame, Notre Dame, IN  46556
%USA  \endaddress
%\email cao.7$\@$nd.edu\endemail
%\address Mei-Chi Shaw,
%Department of Mathematics, University of Notre Dame, Notre Dame, IN  46556
%USA  \endaddress
%\email mei-chi.shaw.1\@nd.edu \endemail
%\address Lihe Wang,
% Department of Mathematics, University of Iowa,
%       Iowa City, IA 52242
%       USA
% \endaddress
%\email lwang\@math.uiowa.edu \endemail
\author   Jianguo Cao,    Mei-Chi
Shaw
     and Lihe Wang**
\endauthor
\thanks{*This paper is a revision of our preprint (May 2003) formerly titled
\lq\lq Estimates for the $\db$-Neumann problem     and nonexistence of
     Levi-flat hypersurfaces  in $\Bbb CP^n$"  where the
nonexistence of $C^{2,\alpha}$ Levi-flat hypersurfaces is proved for
$\alpha>0$.}
\endthanks
\thanks{**All three authors are  partially
supported by NSF grants.   }\endthanks
\address Jianguo Cao,
Department of Mathematics, University of Notre Dame, Notre Dame, IN  46556
USA  \endaddress
\email cao.7$\@$nd.edu\endemail
\address Mei-Chi Shaw,
Department of Mathematics, University of Notre Dame, Notre Dame, IN  46556
USA  \endaddress
\email mei-chi.shaw.1\@nd.edu \endemail
\address Lihe Wang,
     Department of Mathematics, University of Iowa,
          Iowa City, IA 52242 USA
\endaddress
\email lwang\@math.uiowa.edu \endemail














\abstract
Let $\Omega$ be a  pseudoconvex   domain with $C^2$  boundary
in $\Bbb CP^n$, $n\ge 2$.  We prove that the  $\db$-Neumann operator $N$
exists for
     square-integrable forms  on
$\Omega$.
Furthermore, there exists  a number  $\epsilon_0>0$ such that  the operators
     $N$, $\db^*N$, $\db N$ and the Bergman projection
      are regular in the Sobolev space  $W^\epsilon ( {\Omega}) $ for
$\epsilon<\epsilon_0$.
The $\db$-Neumann operator is used to construct $\db$-closed extension on
$\Omega$
for forms on the boundary $b\Omega$.  This  gives solvability for the
tangential
Cauchy-Riemann operators on  the boundary.

Using these results, we  show that there exist no
non-zero $L^2$-holomorphic $(p, 0)$-forms on any   domain
with $C^2$  pseudoconcave boundary in $\Bbb CP^n$ with $p > 0$ and $n\ge 2$.
As a consequence,  we prove   the nonexistence of $C^{2}$
     Levi-flat hypersurfaces  in $\Bbb  CP^n$, a result obtained earlier by Siu
for $C^8$ Levi-flat
hypersurfaces.


\endabstract

\endtopmatter

Preprint is available at  http://www.nd.edu/~meichi/preprints.html


\heading   {\bf 1. Introduction} \endheading

A main result of this paper is the following:
\proclaim{Theorem 1} There exists no $C^{2}$ real Levi-flat
hypersurface   in $\Bbb CP^n$,
$n\ge 2$.
\endproclaim
A $C^2$ real
hypersurface $M$   is called  Levi-flat if its Levi-form vanishes on
$M$.
Theorem 1 is inspired by two
recent papers of Siu [Siu2,3] who proved that there exists no
$C^8$ Levi-flat hypersurface in $\Bbb C P^n$, $n\ge 2$.  The required
smoothness has been reduced to $C^4$ by  Iordan [Io].
       Nonexistence of real analytic
Levi-flat hypersurfaces in $\Bbb CP^n$ was obtained earlier in Lins
Neto [LNe] for
$n\ge 3$.

Our proof of Theorem 1  depends on the $L^2$  existence and
regularity of  the
$\db$-Neumann problem  for   pseudoconvex domains   with $C^2$ boundary
in $\Bbb CP^n$  and is  much simpler then the work cited above.  The
$\db$-Neumann problem
has been studied extensively for pseudoconvex domains in $\Bbb C^n$
or in a Stein manifold. In contrast, the
$\db$-Neumann problem on pseudoconvex domains in $\Bbb CP^n$ has not
been studied systematically before.

    Let $\Omega\subset\subset\Bbb CP^n$
be a pseudoconvex domain with  $C^2$-smooth boundary
$b\Omega$ and let $\delta(x) = d(x, b\Omega)$ be the distance
function from $x\in\Omega$ to
$b\Omega$. We call $t_0 = t_0(\Omega) $   the
order of plurisubharmonicity
for
the distance function  $\delta$ if
$$ t_0(\Omega) = \sup \{ 0<\epsilon \le 1 | i\partial\bar{\partial} (-
\delta^{\epsilon   }) \ge 0 \text{ on } \Omega\}.   $$
    In
$\Bbb CP^n$ with the standard Fubini-Study metric,     Ohsawa-Sibony
[OS] showed  that there
exists $ 0< t_0(\Omega)\le 1  $ for any
pseudoconvex domain $\Omega \subset \Bbb
CP^n$ with $C^2$-smooth  boundary using results of Takeuchi [Taku] and
[DF1]. The other main result in this paper is the following:


\proclaim{Theorem 2} Let $\Omega$ be a pseudoconvex domain   with
$C^2$-smooth boundary in
$ \Bbb CP^n$ and  let  $ t_0$ be the order of plurisubharmonicity
for the distance function $\delta$.
Then the $\db$-Neumann operator
$N_{(p, q)}$    exists on
$L^2_{(p, q) }(\Omega)$ where $0\le p,  q\le n$.   Furthermore,  $N,
\bar{\partial}N,
\bar{\partial}^*N$ and the Bergman projection $   P$ are  exact regular
     on
$W^{s}_{(p, q) }(\Omega)$ for $0 < s <
\frac 12 t_0 $ with respect to the  $W^s(\Omega)$-Sobolev norms.
\endproclaim
    We   give a self-contained
treatment of the $\db$-Neumann problem on pseudoconvex domains in
$\Bbb CP^n$ in Sections 2 and 3.



In Section 4 we use Theorem 2 to study the tangential Cauchy-Riemann
equation on the pseudoconvex
boundary $M=b\Omega$
$$\db_b u= f \quad \text{in } M,\tag 1.1$$
where  $ f$ is a $(p,q)$-form in $M$ satisfying some compatibility
condition with $1\le q\le n-1$.
When $n>2$ and $q<n-1$, the  compatibility condition for equation (1.1)  to
be solvable
is that
$$\db_b f = 0\quad \text{on }M.\tag 1.2$$  When $n\ge 2$ and $q=n-1$, (1.2) is
satisfied trivially and the
compatibility condition for equation (1.1)   is substituted  by  the
following moment
condition:
$$\int_M f\wedge \Psi =0, \tag 1.3$$
where $\Psi$ is  any $\db_b$-closed  $(n,0)$-form  on
$M  $  (see (9.2.12a) of [CS, p216].)
In \cite{Siu2,3},    existence and regularity
of equation (1.1) on a Levi-flat boundary $M$ in $\Bbb CP^n$ is studied.
The two different compatibility
conditions
(1.2) and (1.3) underlie
the reason why  earlier work   differentiate $\Bbb CP^n$ when $n\ge 3$ and
$\Bbb CP^2$.

When $M$ is a $C^2$ Levi-flat  hypersurface in $\Bbb CP^n$,  we   show that the
compatibility condition (1.3) is void    by deriving
several Liouville type theorems which are of independent interest.
In particular, we show that
\proclaim{Proposition  4.5} Let $\Omega^+\subset\subset\Bbb CP^n$ be a
pseudoconcave domain with  $C^2$-smooth boundary $b\Omega^+ $, $n\ge 2$.   Then
$L^2_{(p,0)}( \Omega^+)\cap \text{Ker}(\db)=\{0\}$  for every $0< p\le n$
and  $L^2_{(0,0)}( \Omega^+)\cap \text{Ker}(\db)=\Bbb C$.
\endproclaim
    We use   the dual of the $\db$-Neumann problem (the $L^2$
$\db$-Cauchy problem)
    to prove this
proposition. A weaker version of this proposition was proved in
Henkin-Iordan \cite {HI}.
Theorem 1 follows from this proposition and the result of \cite{OS} easily (see
Section 6 below).

There is a another proof of Theorem 1 for the $C^{2,\alpha}$ case.
We   use the    $W^s(\Omega)$-regularity result for $ N$
         to prove the existence of continuous solutions for equation
$$\partial_b\db_bu=\theta\quad \text{in }M\tag 1.4$$ on a Levi-flat
hypersurface $M$. The
$\partial_b\db_b$ equation is analogous to the classical
$\partial\db$ Poincar\'e-Lelong equation. A  simple geometric proof of
      Theorem 1 for the  nonexistence of $C^{2,\alpha}$ Levi-flat
hypersurface in
$\Bbb CP^n$, $n\ge 2$ and
$\alpha>0$,  along the lines  used by Siu, can be derived from (1.4).
The key
observation is that  the
$\partial_b\db_b$ equation on functions,  when restricted to  each complex
foliation leaf of
$M$, is elliptic. Thus, any solution for (1.4)
has    H\"older regularity on each leaf from the elliptic theory.  To prove
continuity of the solution for equation
(1.4) on a Levi-flat hypersurface,  one only needs to prove  continuity of
the solution in the
transversal direction, which can be proved by using an analogue of
the one-dimensional
Sobolev embedding  theorem.
These arguments  are used in Section 5.
This  approach, which greatly streamlines the original
Siu's arguments,
may have  other applications.


We remark that given $s>0$,  the Bergman projection   $  P$ or the
$\db$-Neumann
operator
$N$ is     not necessarily   bounded
    from  $ W^s_{(p, q)}(\Omega)
     $ to $ W^s_{(p, q)}(\Omega)$ even for  smooth pseudoconvex
domains $\Omega$ in $\Bbb C^n$  (cf. \cite{Ba}, \cite{Chr}).
One can only expect  the boundary regularity of the $\db$-Neumann operator
$N$  to be regular in $W^{s}_{(p, q) }(\Omega)$ with small $s$ for
general pseudoconvex domains in $\Bbb C^n$ or $\Bbb CP^n$, since compact
domains in  $\Bbb C^n$
can
be embedded in $\Bbb CP^n$.



When   $\Omega \subset\subset
\Bbb C^n$ with $C^\infty$-smooth boundary and $\Omega$ has a plurisubharmonic
defining function,   Boas and Straube [BS] prove
that   Theorem 2 holds for any
$s > 0$. For
$\Omega
\subset\subset
\Bbb C^n$ with $C^2$ boundary,  Berndtsson-Charpentier [BC] (see also Kohn
\cite{Ko2}) obtain the
Sobolev regularity for
the operators $ \bar{\partial}^*N$ and the Bergman projection $  P$.
Boas-Straube [BS1] show  that
regularity for the
Bergman projection and the $\db$-Neumann operator are equivalent for
$C^\infty$-smooth pseudoconvex domains in
$\Bbb C^n$ (and also  for smooth pseudoconvex domains in complex manifolds with
strongly plurisubharmonic   functions in a
neighborhood of the boundary).  For
     pseudoconvex domains $  \Omega \subset \Bbb CP^n $,  such a
smooth  strictly plurisubharmonic function $\phi$ might not exist in a
neighborhood of the boundary
$b \Omega$. Therefore, the proof of   Theorem 2 has extra difficulties
when compared
to the case of domains
in  $\Bbb C^n$.  We use results by  Takeuchi [Taku] and
        Ohsawa-Sibony [OS]   for
pseudoconvex domains  with $C^2$-smooth  boundary  in $ \Bbb
CP^n$.
    We remark that our proof of the regularity for the Bergman
projection and $\db^*N$ is similar to the proof
in \cite{BC}, but our proof of
the regularity results for  the $\db$-Neumann operator $N$ and $\db N$ in
Theorem 2
are new, even for   domains with
$C^2$-smooth boundaries in
$\Bbb C^n$.




It should be pointed out that there exist   nonsmooth  Levi-flat real
hypersurfaces in
$ M\subset \Bbb CP^2$  in the sense that   $\Bbb CP^2\setminus M$
consists of two pseudoconvex
domains.  For example, let $$ M=\{[z_0,z_1,z_2]\in\Bbb CP^2\mid
|z_0|=|z_1|\},$$
where $[z_0,z_1,z_2]$ are homogeneous coordinates in $\Bbb CP^2$.
Then $M$ is Levi-flat since $\Bbb CP^2\setminus M$ consists of two
pseudoconvex domains. The hypersurface $M$
is   smooth except at $[0,0,1]$, where $M$ is not $C^1$.
In addition, the positive curvature condition
is needed for the possible  generalization of Theorem 1. One can inspect
the example
$S^1 \times \Bbb CP^1$, which is a $C^\infty$-smooth
Levi-flat hypersurface in $\Bbb CP^1 \times \Bbb CP^1$. Clearly,
$\Bbb CP^1 \times \Bbb CP^1$ does not  admit a K\"ahler metric with
positive holomorphic bisectional curvature, (e.g., cf. [HSW]).
Any compact K\"ahler manifold with positive  holomorphic bisectional curvature
must be biholomorphic to $\Bbb CP^n$ (e.g., cf. [Mok]). Although results in
this paper
are stated for domains in $\Bbb CP^n$, the conclusions of these results
remain to be true for corresponding domains  with positive holomorphic
bisectional curvature as
well.







       The plan of this paper is as
follows:  In Section 2 we prove
$L^2$ existence theorems of $\db$ and  the
$\db$-Neumann problem on pseudoconvex domains with $C^2$ boundary in
$\Bbb CP^n$. In
Section 3 we prove that the $\db$-Neumann operator on pseudoconvex
domains with
$C^2$ boundary is regular in the Sobolev spaces as stated in Theorem 2.
Using the
existence and regularity of the
$\db$-Neumann operator,  one can study  the $L^2$ solution of the
$\db$-equation  with prescribed support on pseudoconvex domains,  including the
case of forms with top degrees. Such $L^2$ Cauchy problem is the key
ingredient in our proof
of Theorem 1.  We
construct $\db$-closed   extensions
from the boundary and prove the existence and regularity of equation (1.1)
in   Section 4. The Liouville type results for holomorphic forms are
also derived in Section 4.
    In Section 5 we study the $\partial_b\db_b$ equations on Levi-flat
hypersurfaces and  prove
Theorem 1 for
$C^{2,\alpha}$ case  when
$\alpha>0$. Theorem 1 is proved in Section 6 using Proposition 4.5.
In the appendix we
include  some    results  for hypersurfaces   in $\Bbb CP^n$ in
differential geometry and some other
    useful facts which are used in this paper. These results in the appendix
are not new
but are included here   for the benefit
of the reader.
\medskip
\noindent
     {\bf Acknowledgment:} The authors would like to thank    A. Iordan and
S. Vassiliadou
for their comments on the original manuscript. In the earlier manuscript,
Theorem 1 was
proved   for
$C^{2,\alpha}$ Levi-flat hypersurfaces only.  The authors   are
especially grateful to Professor  T. Ohsawa who kindly pointed out at
the Hayama conference,
December 2003, in Japan that  one can prove   Theorem 1  directly
from Proposition 4.5, which led to  the improved results. The proof
of Theorem 1  given
in Section 6 follows from his suggestions.



\end






















